%! TEX root = ../bcd-19-20.tex

\chapter{Examen 2018--2019}

\section{Numărul 1}

\indent\indent 1. Să se arate că funcția:
\[
  f(x, y) = x \cdot \ln \frac{x - y}{x + y} + y\cdot e^{\frac{x}{y}}
\]
verifică ecuația:
\[
  x \frac{\partial f}{\partial x} + y \frac{\partial f}{\partial y} = f(x, y).
\]

\vspace{1cm}

2. Să se determine extremele funcției $ y = y(x) $ definită implicit de
ecuația:
\[
    x^3 + y^3 - 3x^2y - 3 = 0.
\]

\vspace{1cm}

3. Să se studieze continuitatea și existența derivatelor parțiale în $(0, 0) $
pentru funcția:
\[
  f(x, y) = \begin{cases}
    \dfrac{xy^2}{x^2 + y^4}, & (x, y) \neq (0, 0) \\
    0, & (x, y) = (0, 0)
  \end{cases}
\]

\vspace{1cm}

4. Să se determine mulțimea de convergență a seriei de puteri
\[
  \sum_{n \geq 1} \dfrac{n+1}{n!} x^n
\]
și apoi să se afle suma seriei numerice
\[
  \sum_{n \geq 1} \dfrac{30^n(n + 1)}{n!}.
\]

\vspace{1cm}

5. Să se dezvolte în jurul lui $ 0 $ funcția:
\[
  f(x) = \dfrac{1}{x^2 - 7x + 10},
\]
precizînd și mulțimea pe care e valabilă dezvoltarea.


%%%%%%%%%%%%%%%%%%%%%%%%%%%%%%%%%%%%%%%%%%%%%%%%%%%%%%%%%%%%%%%%%%%%%%
\section{Numărul 2}

\indent\indent 1. Fie funcția:
\[
  f(x, y) = (x^2 - y^2) \arctan \dfrac{x}{y} + (x^2 + y^2) \sin \dfrac{x}{y}.
\]
Arătați că e satisfăcută relația:
\[
  x \dfrac{\partial f}{\partial x} + y \dfrac{\partial f}{\partial y} = 2 f(x, y).
\]

\vspace{1cm}

2. Să se determine extremele locale ale funcției:
\[
  f : \RR^2 \to \RR, \quad f(x, y) = xy^2 e^{x-y}.
\]

\vspace{1cm}

3. Să se studieze continuitatea și derivatele parțiale în origine pentru funcția:
\[
  f(x, y) = %
  \begin{cases}
    \dfrac{x^3 - y}{\sqrt{x^2 + y^2}}, \quad (x, y) \neq (0, 0) \\
    0, & (x, y) = (0, 0)
  \end{cases}
\]

\vspace{1cm}

4. Determinați mulțimea de convergență a seriei $ \ds\sum_{n \geq 1} \dfrac{n}{n+1}x^n $
și apoi suma seriei $ \ds\sum_{n \geq 1} \dfrac{n}{(n+1)2019^n} $.

\vspace{1cm}

5. Să se dezvolte în serie de puteri ale lui $ x $ funcția $ f(x) = \ln(1 + 2x) $
și precizați mulțimea pe care e valabilă dezvoltarea.


%%%%%%%%%%%%%%%%%%%%%%%%%%%%%%%%%%%%%%%%%%%%%%%%%%%%%%%%%%%%%%%%%%%%%%
\section{Restanță}

\subsection{Numărul 1}

\indent\indent 1. (a) Studiați convergența seriei:
\[
  \sum_{n \geq 1} \dfrac{a^n}{n^\alpha}, \quad a > 0, \alpha \in \RR.
\]

(b) Studiați convergența absolută a seriei:
\[
  \sum_{n \geq 1} \dfrac{(-1)^n}{2^{2n} \cdot n!}
\]

\vspace{1cm}

2. Studiați convergența simplă și uniformă pentru șirul de funcții:
\[
  f_n : [2, \infty) \to \RR, \quad f_n(x) = \dfrac{n^2 x}{n^2x^2 + 1}.
\]

\vspace{1cm}

3. Găsiți extremele globale ale funcției:
\[
  f: \RR^3 \to \RR, \quad f(x, y, z) = x + 2y - 2z
\]
pe mulțimea $ K = \{(x, y, z) \in \RR \mid x^2 + y^2 + z^2 \leq 16 \} $.

\vspace{1cm}

4. Arătați că funcția:
\[
  f(x, y) = (x^2 - y^2) \cdot \ln \dfrac{x - y}{x + y}
\]
verifică ecuația:
\[
  x \dfrac{\partial f}{\partial x} + y \dfrac{\partial f}{\partial y} = 2f(x, y).
\]

\vspace{1cm}

5. Să se dezvolte în jurul originii funcția $ f(x) = \dfrac{1}{x^2 + 5x + 6} $
și să se precizeze intervalul pe care această dezvoltare este valabilă.

%%%%%%%%%%%%%%%%%%%%%%%%%%%%%%%%%%%%%%%%%%%%%%%%%%%%%%%%%%%%%%%%%%%%%%
\subsection{Numărul 2}

\indent\indent 1. (a) Studiați convergența seriei 
$ \ds\sum_{n \geq 1} \dfrac{n \sqrt{n}}{\left( a + \dfrac{1}{n} \right)^n} $,
unde $ a > 0 $.

(b) Studiați convergența absolută pentru seria $ \ds\sum_{n \geq 1} (-1)^n \dfrac{(2n)!}{n^n} $.

\vspace{1cm}

2. Studiați convergența simplă și uniformă pentru șirul de funcții:
\[
  f_n : [1, \infty) \to \RR, \quad f_n(x) = \dfrac{n}{nx + x^3}.
\]

\vspace{1cm}

3. Găsiți extremele funcției:
\[
  f: \RR^3 \to \RR, \quad f(x, y, z) = -4x - 3y + 6z
\]
pe mulțimea:
\[
  K = \{(x, y, z) \in \RR^3 \mid x^2 + y^2 + z^2 \leq 1 \}.
\]

\vspace{1cm}

4. Arătați că funcția:
\[
  f(x, y) = xy + x \cdot \sin\dfrac{y}{x}
\]
verifică ecuația:
\[
  x \dfrac{\partial f}{\partial x} + y \dfrac{\partial f}{\partial y} = xy + f(x, y).
\]

\vspace{1cm}

5. Dezvoltați funcția $ f(x) = \dfrac{1}{x^2 + 4x + 3} $ în serie de puteri în jurul
originii și precizați intervalul pe care dezvoltarea este valabilă.

